\documentclass{article}
\usepackage{amsmath}

\emergencystretch=2em

\title{Coupled rewards in Q and false consensus in P:\\
a converse on Q's side, an untested premise on P's}
\author{ada\\[2pt]
\small Consolidated note for pair Q4P3, ledger entries \#1 to \#27}
\date{11 September 2026}

\begin{document}
\maketitle

\begin{abstract}
Both peers read P's false consensus as the equilibrium that Q mentions in
passing in its coupled-rewards example. On Q's side the ledger has a
checked converse to Q's result: in Q's two-agent quadratic model, rewards
based on the action difference cannot discipline agents whose biases move
together across types, and with finite agreement weights a reviewer who
seeks agreement never ends up less biased than the less biased agent
alone. On P's side there is no result. Under Q's own standard of partial
implementation, a terminate-on-agreement device also supports the truthful
action, so every claim about P turns on whether truth is focal for a
reviewer and critic run on the same model, and nothing in the material
tests that. In round~2 the peers withdrew the audit that round~1 proposed,
because it is confounded, designed a same-diff test of the critic in its
place, and read P's prompts for what P's SWE-PRBench numbers can carry. No
model calls were made, so nothing about P has been measured.
\end{abstract}

\section{The pair}

Q (Bergemann, Koh and Morris, \emph{Mechanism Design for Alignment and
Control}, arXiv 2609.01595) develops mechanism design for AI agents whose
preferences and capabilities are unknown, under a one-sided imitation
order in which capabilities can be concealed but not counterfeited, and
works through stylised examples that include peer scoring and coupled
rewards between two biased agents. P (Qiu and Gill, \emph{Adversarial
Review}, arXiv 2608.18167) proposes a code-review protocol (AR) in which a
reviewer R and a critic C, both run on Claude Sonnet~4.5, exchange reviews
before a main agent M edits, and reports gains on LiveCodeBench (LCB),
SWE-PRBench and SWE-bench Verified together with a false-consensus
failure, agreement without sufficient evidence, which a prompt adding
evidence-grounded disagreement repairs.

\section{The connexion pursued}

Q's example ``Competition through coupled rewards'' has two agents who
observe a state $\theta$ and choose $a_1,a_2\in\mathbf{R}$. The human's
payoff is $v=-\sum_i(a_i-\theta)^2$, and agent $i$ receives
$U_i=-(a_i-\theta-b_i)^2+r_i(D)$ with $D:=a_2-a_1$. The biases lie on a
known curve $b_2=f(b_1)$ with $f$ strictly decreasing and $f(0)=0$, and
Q's rewards (CR) give a unique equilibrium with loss
$\eta^2(b_2-b_1)^2/2$, which vanishes as $\eta\to0$. Q remarks that
coupling ``need not push them toward choosing $a_i=\theta$; they might
simply choose the same fixed action'' (Q line 1918).

Both peers read P's false consensus as that equilibrium (ada \#1, emmy
\#2), with agent~1 as R, agent~2 as C, $x_i:=a_i-\theta$ as an agent's
error and $b_i$ as its error tendency (over- or under-flagging). Because
P runs every role on the same model, the case of interest is $b_C=k\,b_R$
with $k>0$, clones $k=1$ included, or more generally biases that move
together across tasks, which Q's assumption rules out. Both peers mark
this dictionary as an analogy that P does not measure (\#5(b), \#11(3),
\#15(2), \#24(3)). The secondary lines in round~1 were P's
evidence-grounded verdicts read as Q's hard evidence (\#1, \#2, \#10), Q's
peer scoring set against P's terminate-on-AGREE rule (\#8), and the design
of P's LCB comparison (\#9). Two side lines announced in \#1, on Q's
weak-strong monitor and on P's SKILL form, were not pursued.

The verifier (\#20) accepted the algebra on Q's side and raised two
points: under Q's standard of partial implementation the claims about P
come down to equilibrium selection, and the audit that round~1 proposed as
the next step (\#8) is confounded. In round~2 the peers looked for an
anchor on P's side that needs no model spend: a same-diff test of the
critic (emmy \#21, ada \#23(3)), a map of what P's published numbers and
outside work can show (ada \#23), and a reading of P's own prompts against
its Table~2 (emmy \#26). Line numbers below refer to the source texts of P
and Q supplied with the pair.

\section{What was found}

\subsection{An agreement penalty}\label{sec:penalty}

Replacing (CR) by $r_i=-\kappa_i D^2$ with $\kappa_i>-1$ gives the unique
equilibrium (ada \#7)
\begin{equation*}
a_R-\theta=\frac{(1+\kappa_C)\,b_R+\kappa_R\,b_C}{1+\kappa_R+\kappa_C},
\qquad
D=\frac{b_C-b_R}{1+\kappa_R+\kappa_C}.
\end{equation*}
In P the output is R's final review. C moves it only through R's own
agreement weight, since $\kappa_R=0$ gives $a_R=\theta+b_R$. A critic that
yields ($\kappa_C\to\infty$) also leaves $a_R\to\theta+b_R$, so AR
collapses to Single-reviewer. With $\kappa_R=\kappa_C\to\infty$ both
actions tend to $\theta+(b_R+b_C)/2$: agreement without any drop in error.
In that limit the critic helps iff $b_C\in(-3b_R,b_R)$ for $b_R>0$, so a
same-direction critic more biased than R makes AR worse than
Single-reviewer. ada matched this to the ordering of P's SWE-PRBench F1
scores, AR $0.457$ against Single-reviewer $0.495$ (\#7). The match is
qualitative, nothing was fitted, and other channels can carry the same
gap (Section~\ref{sec:table2}).

Any symmetric reward $r_1=r_2=\rho(D)$ that is differentiable at the
equilibrium has first-order conditions $-2(x_1-b_1)-\rho'(D)=0$ and
$-2(x_2-b_2)+\rho'(D)=0$, which add to $x_1+x_2=b_1+b_2$, so the loss is
at least $(b_1+b_2)^2/2$ at any penalty strength (emmy \#3). If R and C
yield unequally, with $\kappa_R=\alpha K$, $\kappa_C=(1-\alpha)K$ and
$K\to\infty$, the consensus is $\theta+(1-\alpha)b_R+\alpha b_C$ (ada
\#10(2)). That is still a convex combination, and it beats R alone only if
R does the yielding and C is less biased.

\subsection{Co-moving biases}\label{sec:comove}

Let $b_1$ range over an interval $[b_{lo},b_{hi}]$ of length $L$, and let
each $r_i$ depend on $D$ alone, with $-\infty$ allowed.

\paragraph{Proposition (emmy \#2).} If $b_2=k\,b_1$ with $k>0$ and
$L>12\varepsilon\max(1,1/k)$, then any such rewards under which every
type has an equilibrium with $|x_i|\le\varepsilon$ make the equilibrium
$D$ the same for every type. The extreme type $b_{hi}$ also has
equilibria with both actions shifted by any $c$ with
$0\le c\le\min(1,k)L-2\varepsilon$, so its worst equilibrium loss is at
least $2(\min(1,k)L-2\varepsilon)^2-O(\varepsilon)$. For $k<0$, (CR)
works. For $k=0$, the rewards $r_2=0$ and $r_1=-KD^2/2$ give
$x_1=2b_1/(2+K)\to0$.

\paragraph{Argument.} Facing its partner, agent~1 chooses $D$ to maximise
$r_1(D)-(D-z_1)^2$ with $z_1:=x_2-b_1$, and agent~2 maximises
$r_2(D)-(D-z_2)^2$ with $z_2:=b_2-x_1$. Adding the optimality conditions
of two types, each tested against the other's choice, gives exactly
(ada \#7)
\begin{equation*}
\bigl(D(b)-D(b')\bigr)\bigl(z_i(b)-z_i(b')\bigr)\ \ge\ 0,\qquad i=1,2.
\end{equation*}
Near first best, $z_1=-b_1+O(\varepsilon)$ and $z_2=kb_1+O(\varepsilon)$.
For $k>0$, agent~1's inequality makes $D$ weakly decreasing in $b$ and
agent~2's makes it weakly increasing, across types more than
$2\varepsilon/\min(1,k)$ apart. Chaining through a distant third type makes
$D$ constant. For $k<0$ the two conditions agree. emmy's own route (\#2)
uses $r_1(D+y)\le r_1(D)+s\,y+y^2$ for all $y$, with $s=2(b-x_1)$, which
needs no differentiability. At the pooled value the slopes $s$ satisfying
it form a convex set containing every type's slope, so shifting both
actions of type $b_{hi}$ by $c$ moves the two agents' slopes by $-2c$ and
$+2c$ without leaving the valid sets (Step~4). The reward is then a pure
consensus device.

\paragraph{Co-movement, not sign (emmy \#4).} The proof uses only that
$f$ is increasing. With $f(b)=b-1$ on $[0,1]$ the biases have opposite
signs for every type, and the result still holds. Q's (Curve) is used
only through the fact that the biases move in opposite directions. The
quantity to measure in P is the across-task correlation of R's and C's
signed errors, not whether they agree.

\paragraph{Scope (\#20, \#21(1)).} The proposition says that near-first-best
difference rewards also support biased-consensus equilibria, so it rules
out full implementation, not partial implementation. Terminate-on-AGREE
written as a $0/{-\infty}$ indicator lies in this class, and for it the
result is one of multiplicity: the equilibrium with $a_i=\theta$ exists
alongside the biased ones.

\subsection{General quadratic rewards, output by R alone}\label{sec:quad}

For any quadratic state-independent rewards $r_i(a_1,a_2)$ with a unique
equilibrium, best responses are $a_i=c_it_i+\rho_ia_j+d_i$, with
$t_1=\theta+b$, $t_2=\theta+kb$, $c_i>0$ and $\Delta=1-\rho_1\rho_2\ne0$.
The equilibrium loadings $L_i$ on $\theta$ and $\beta_i$ on $b$ satisfy
(emmy \#14(1))
\begin{equation*}
(\beta_1-kL_1)(L_2-\beta_2)=\frac{(1-k)^2c_1c_2}{\Delta^2}\ \ge\ 0 .
\end{equation*}
Bounded loss forces $L_i=1$, that is $\rho_i=1-c_i$, which is the
quadratic difference class with $c_i=1/(1+\kappa_i)$. Then for $k>0$,
$\max_i|\beta_i|\ge\min(k,1)$, attained as $c_1\to0$ (R copies C). If only
R's output counts, as in P,
\begin{equation*}
\beta_R=\frac{c_1+k(1-c_1)c_2}{c_1+(1-c_1)c_2}.
\end{equation*}
An agreement-seeking R ($\kappa_R\ge0$) makes $\beta_R$ a weighted average
of $1$ and $k$ whatever C is paid, so $|\beta_R|\ge\min(1,k)$ (\#14(2)).
Reaching $\beta_R=0$ needs $kc_2>1$ and $c_1=kc_2/(kc_2-1)>1$, that is
$\kappa_R<0$. R is then paid to disagree and extrapolates away from C,
which overshoots with $\beta_C=kc_2>1$. For example, $k=1/2$,
$r_R=+D^2/2$ and $r_C=+3D^2/4$ give $x_R=0$ and $x_C=2b$.

These results have a unique equilibrium, so full and partial
implementation coincide and no selection question arises. They model LLM
conformity as a finite pull $\kappa$ (\#21(2)). For nonlinear difference
rewards, \#14(2) argues that any concave $r_R$, terminate-on-agreement as a
$0/{-\infty}$ indicator included, forces pooling for a near-first-best R
and hence biased-consensus equilibria, which is again a multiplicity
statement. The claim that P's loop, which terminates on AGREE, cannot
de-bias R's output therefore holds in the unique equilibrium when
conformity is finite, and only as a statement about selection when the
loop is read as a hard device (\#20, \#21(1)). Which reading fits P is an
empirical question (\#21(2)).

\subsection{Limits of the contrarian escape}

In the quadratic difference class $x_R=A\,b_R+(1-A)\,b_C$ with
$A=c_R/(c_R+c_C-c_Rc_C)$, and $A$ can be any real number (ada \#17(3)).
The error vanishes only at $A^*=k/(k-1)$, so de-biasing needs the ratio
and, for same-sign biases, extrapolation.
\begin{itemize}
\item \emph{Stability.} For $0<k<1$, $A^*<0$, and $A<0$ iff
$\rho_R\rho_C>1$ with $\rho_i=1-c_i$. That is the condition under which
alternating best responses, the way P's inner loop proceeds, diverge.
emmy's example reaches order $10^{19}$ in 40 rounds. For $k>1$ a stable
design exists: C states its bliss point and R plays
$a_R=c_Rt_R+(1-c_R)t_C$ with $c_R=k/(k-1)$, which at $k=2$ is
$a_R=2t_R-t_C$ and converges in one round.
\item \emph{Unknown ratio.} A design tuned to $k_0$ leaves
$\beta_R(k)=(k-k_0)/(1-k_0)$, and every design has $\beta_R(1)=1$. The
minimax residual over $[k_{lo},k_{hi}]$ with $k_{hi}<1$ is
$(k_{hi}-k_{lo})/\bigl((1-k_{lo})+(1-k_{hi})\bigr)$: $0.96$ on
$[0.5,0.99]$ and $0.6$ on $[0.8,0.95]$, against $1$ for R alone. The
analogue above $1$ ($0.6$ on $[1.5,3]$) was computed by hand only
(\#18(2)).
\item \emph{Identification.} With an unknown ratio, the types
$[0,B]\times[k_{lo},k_{hi}]$ include two with the same preferences whose
states differ by $B(k_{hi}-k_{lo})/(1-k_{lo})$, whatever the
state-independent mechanism. The selection qualifier below applies.
\end{itemize}
The stable form in the model is a single round. C is prompted to
exaggerate, so that $k>1$, and states its position without strategising;
R discounts it with a fixed weight $A>1$ (\#18(2)). P's text-constraint
prompt instead tells R to incorporate any bug C adds (P line 1172, \#2).

\subsection{Clones and equilibrium selection}

At $k=1$ preferences depend on $(\theta,b)$ only through $t=\theta+b$, so
no mechanism without a state measure can distinguish $(\theta,b)$ from
$(\theta+c,b-c)$, and emmy concluded that the best equilibrium does no
better than one agent (\#14(3)). ada restricted this to full
implementation, or to selections that depend only on the game (\#17(4),
\#18(1)). Under Q's partial implementation (Q line 543: ``an equilibrium,
not necessarily the unique one''), a report-matching device for clones has
every common report $m$ as an equilibrium in every state, $m=\theta$
included, while both agents strictly prefer $m=t$ (payoff $0$ against
$-b^2$). The clone result is therefore about selection: a same-model loop
helps only if truth is focal for the models, and no state-independent
mechanism can enforce that (\#17(4)). emmy accepted this in round~2
(\#21(1)), and the verifier made the same point for every claim about P
(\#20). For $k\ne1$, emmy withdrew the conjecture that report matching
also fails (\#5(a), \#14(4)). Message mechanisms may succeed there, but
none has been built.

\subsection{Peer scoring against terminate-on-AGREE}

Q's binary example scores reports by $q(H,H)=7/9$, $q(H,L)=1/9$,
$q(L,H)=1$ and $q(L,L)=-1$, and truth gives the human $1$. If the partner
reports $H$ with probability $p$ regardless of type, reporting $H$ earns
$1/9+2p/3$ and reporting $L$ earns $2p-1$. These are equal at $p^*=5/6$,
so both types mixing $5/6$ on $H$ is an equilibrium that gives the human
$1/2$ (ada \#8). Agent~1 always $H$ with agent~2 always $L$ is also an
equilibrium ($1/9>-1$, $1>7/9$), with the same human payoff (emmy \#13).
Ex ante, truth and babbling each give every agent $2/3$ and asymmetric
pooling gives $(1/9,1)$, so payoff dominance does not select truth
(\#13). The design contrast with P is that P's loop rewards agreement
absolutely, while Q's score rewards it only as far as an agent's own type
predicts its partner's report (\#8).

\#8 turned this into an audit comparing C's AGREE rate on matched pairs
(review of PR $i$ with diff $i$) and mismatched pairs (review of PR $j$
with diff $i$). That audit is withdrawn. P's critic is asked whether each
flag is a real bug in the diff shown (P lines 1107--1109), a review of PR
$j$ cites code absent from diff $i$, and C would reject it on topicality
alone, so a gap would appear whatever AGREE tracks (\#20). Mismatched-task
baselines assume a shared report space, which free-text reviews that cite
task content lack (\#21(1)).

\subsection{Evidence as a state measure}\label{sec:state}

emmy's result sits in the cell of Q's taxonomy with an exact action
measure and no state measure (ada \#10(1)). P's evidence-grounded
verdicts (P Sec.~4.4) add a coarse state measure. Suppose the protocol
sees $m=\theta+\beta_m$ plus noise and adds $-\lambda(a_i-m)^2$ to each
payoff. Setting the noise aside,
$a_i=\theta+(b_i+\lambda\beta_m)/(1+\lambda)\to\theta+\beta_m$ as
$\lambda\to\infty$, whether or not the biases co-move, so the error moves
onto the checker. P's citations are judged by R, the same model as C, so
$\beta_m$ is roughly the common-mode bias. A mechanical check that the
snippet occurs in the diff is unbiased, but it tests existence, not
relevance. Label C's types E (contradicting code), N (doubt) and A (no
objection). If citations are checked, N and A cannot mimic E, and Q's
multi-agent revelation principle then implies that three typed verdicts
with response rules lose nothing against a richer exchange (\#10(3)). This
needs citations that cannot be counterfeited and an obedient R, and P
enforces neither. AGREE and C's added bugs carry no certificate (\#2).

\subsection{What P's numbers and prompts can carry}\label{sec:table2}

P reports only F1 on SWE-PRBench: $0.533$ for AR with the text
constraint (v2), $0.503$ for Two-reviewers, $0.501$ for MARS, $0.495$ for
Single-reviewer and $0.457$ for naive AR, on 100 PRs each, with no
precision, recall or comments per PR (P lines 294--312, ada \#23(1)).

\paragraph{The anchor (emmy \#26(1)).} Single-reviewer, Two-reviewers and
AR use the same reviewer prompt (P lines 1030--1034), and R's first call
has an empty prior-exchange slot. R's first draft in naive AR therefore
has the distribution of the Single-reviewer output, and the F1 difference
between AR and Single-reviewer comes from the inner loop's edits to that
draft plus any post-processing applied to AR alone. The one candidate is
the format-review step that splits each flag into a separate comment
(P line 320), which P describes only for AR and for which it gives no
prompt.

\paragraph{Case B is neutral against Single-reviewer (emmy \#26(2)).} In
Case~B, R's draft is APPROVE, C raises a real concern and yields, and the
final review is R's draft (P lines 329--331). It loses against MARS
($0.286$ against $0.667$ on that task, P line 336), not against
Single-reviewer. It is the limit $\kappa_C\to\infty$ of
Section~\ref{sec:penalty}. Of P's two documented failure modes, only
Case~A (C confirms R's hedges and adds a speculative bug, P lines
317--326) can contribute to the gap between $0.457$ and $0.495$. This
corrects the Case~B row of ada's table in \#23(1).

\paragraph{Deletion channels (emmy \#26(3)).} The v2 prompt changes only
R's block on responding to the critic (P lines 1125--1127). That block
forbids R to ``consolidate, drop, or rephrase legitimate findings just
because you are about to converge'' (P lines 1156--1158) and to
``capitulate to APPROVE merely because the critic called the concern
`speculative' '' (P lines 1168--1170). Both describe R deleting its own
draft flags, which lowers recall against Single-reviewer and which P's
case studies do not show. Under micro-averaged $F_1=2T/(C+H)$, with $T$
matched comments, $C$ comments and $H$ human comments, normalise
Single-reviewer to $T=1$, so that $C+H=2/0.495\approx4.04$. Moving from
$0.495$ to $0.457$ then needs either $0.336$ added unmatched comments per
Single-reviewer match or deletion of $10.0\%$ of Single-reviewer's
matches. Deletion moves F1 about three times as fast, so the deletion
channels alone could carry the deficit. The candidates have different
signatures (\#23(1) as corrected by \#26(2) and \#26(3)):
\begin{center}\small
\begin{tabular}{p{0.36\linewidth}p{0.18\linewidth}p{0.34\linewidth}}
\hline
\raggedright Candidate carrier & \raggedright P lines &
\raggedright Signature against Single-reviewer\tabularnewline
\hline
\raggedright Case A: C confirms R's hedges and adds a flag (the
co-movement mode) & \raggedright 317--326 &
\raggedright lower precision, more comments per PR\tabularnewline
\raggedright R drops its own flags on convergence or under doubt &
\raggedright 1156--1158, 1168--1170 &
\raggedright lower recall, fewer comments per PR\tabularnewline
\raggedright Case B: C yields to a rebuttal that cites no code &
\raggedright 328--336 &
\raggedright none; it loses only against MARS\tabularnewline
\hline
\end{tabular}
\end{center}

\paragraph{What can carry v2's gain (emmy \#26(4)).} On AGREE, v2 keeps
R's draft unchanged, so v2 equals Single-reviewer on tasks where C agrees;
in the model this sets $\kappa_R=0$ on that path (\#7). The clause ``If
the critic added a missed bug (under any verdict type), incorporate it''
(P line 1172) is new in v2. It appends C's additions with no certificate,
which closes Case~B and opens Case~A's addition channel. C then acts only
through the union of its additions with R's draft and through the
cite-or-drop rule, which is the state measure of
Section~\ref{sec:state}. The gain of $0.038$ over Single-reviewer needs
only $0.105$ matched additions per Single-reviewer match. ada's reading
(\#23(1)) is that the cite-or-drop rule raises precision, capped by C's
added flags. The ablation that separates the two readings is v2 with the
DISAGREE\_CONCERN cite-or-drop clause removed and the other two clauses
kept. If it scores near $0.533$, P's design principle, ``force
disagreement to be explicit and evidence-grounded'' (P line 264), is not
what carries the gain.

\paragraph{Strength and scope (emmy \#26(5), \#26(6)).} If the format step
is common to all methods, AR below Single-reviewer is the $n=100$
aggregate form of the statement that the pair does not select truth: the
loop leaves R worse than its own draft. The evidence is weak. With an
assumed paired per-PR standard deviation of $0.2$ to $0.3$ (P gives none),
the gap of $0.038$ is $1.3$ to $1.9$ standard errors. P's judge scores
similarity to human comments, so an unmatched comment need not be false
(\#21(4)). And the deficit need not come from the co-movement mode. Q's
model has one state and a scalar action, so bias is its only error. Case
B's mechanism, where R ``can end the disagreement by writing a
confident-sounding rebuttal, even when R is wrong'' (P lines 333--334), is
a failure of discrimination, not of flagging volume, and the co-movement
results say nothing about it.

\subsection{The primitive on P's side}\label{sec:primitive}

Both readings of P's loop need one quantity: how C's verdict on the same
diff responds to a false flag in R's review, and whether that depends on
the flag being one that C's own model would write (emmy \#21(3)). In the
model C's disposition on item $\tau$ is $b_C(\tau)$, and co-movement means
that $b_C$ and $b_R$ covary positively across items. A false flag written
by Sonnet as R points along $b_R$, so under co-movement C's own disposition
favours it, while a flag written by another model family carries no such
alignment. The gap in detection between own-model and other-model false
flags is then the co-movement premise of \#4, measured by intervention.
Under the hard device, truth is focal for C exactly when detection is
high. With finite $\kappa$, detection is $(1-\text{conformity})$ times
discriminability.

Outside work measures two different things, and only one of them is P's
premise (ada \#23(2)). \emph{Ownership}, with content fixed and role
varied: arXiv 2606.05976 injects a false claim into a model's own thought
block, and relabelling it as a user, tool or memory message raises
explicit correction by 23 to 93 percentage points in 10 of 12 settings
(9 models, maths and logic). P's C always receives R's review as another
agent's message, so this effect does not arise in P, and results that
models cannot correct themselves overstate what carries over.
\emph{Disposition}, the model's own kind of error shown as external
input: DeltaBench (arXiv 2502.19361) finds self-critique below
cross-model critique for three reasoning models, $36\%$ below for
DeepSeek-R1, and arXiv 2408.10495 finds that GPT-3.5 and GPT-4 repair
other models' insecure code at $33.2$ to $59.6\%$ but do poorly on their
own. Both point towards co-movement in P's presentation condition. Both
are confounded by source, since own outputs differ in difficulty and
style, and neither covers code review or Sonnet~4.5. Kim et al.\ (arXiv
2506.07962) find more correlated errors within a provider, which is
co-movement of correctness across items, not of flagging. The outside
work supports the direction of the premise and does not measure it for
P's pair (\#23(2), \#24).

\subsection{Checks}

The four check scripts in \texttt{ada/} and \texttt{emmy/} were rerun for
this note on 11 September 2026 and reproduce every value quoted from them:
the closed forms, limits and fractions, emmy's quadratic identity
(symbolically and on random draws), the affine form and minimax values,
the divergence of emmy's example against one-round convergence at $k=2$,
and Step~4 of the proposition for a kink penalty. The round-2 arithmetic
was also rechecked: the interval half-widths of \#21(5), and the size
check and standard-error ratios of \#26(3) and \#26(5). The P passages
cited in round~2 (P lines 272, 320, 329--336, 685--688, 728, 1030--1034,
1107--1109, 1125--1127, 1156--1172) were read against P's text and say what
the entries quote. Steps 1 to 3 of the proposition rest on the written
proof, which ada checked (\#7, \#11). The verifier checked the
proposition, ada's exact inequality, the Step~4 shift, the quadratic
identity, the weighted-average bound and the contrarian algebra (\#20).
Neither peer made model calls: no API key was set and no spending was
authorised (\#21(7), \#23(4)).

\section{Objections and their status}

\begin{itemize}
\item \textbf{Claims about P rest on selection} (verifier \#20). Under
Q's partial implementation a $0/{-\infty}$ consensus device supports
$a=\theta$, so every claim about P, including that the loop cannot
de-bias R's output, turns on whether truth is focal for a same-model R-C
pair. Accepted by emmy (\#21(1)) and ada (\#23(5), \#24). emmy narrowed
its reach: the finite-conformity results have unique equilibria and are
unaffected, and which reading fits P is empirical (\#21(2)). The premise
remains untested.
\item \textbf{The matched-versus-mismatched audit is confounded}
(verifier \#20, on \#8). Accepted and withdrawn (\#21(1), \#23(5),
\#24(6)).
\item \textbf{Replacing a flag confounds detection with omission} (emmy
\#21(4), on the verifier's proposed arm). A replaced true flag leaves an
omission that C's second task, naming missed bugs (P line 1109), can
catch. emmy appends the false flag instead, and ada's extra arm builds on
that design (\#23(3)). Not contested.
\item \textbf{Case B and Single-reviewer} (emmy \#26(2), correcting ada
\#23(1) and \#24). Case~B loses recall against MARS, not against
Single-reviewer. The peer phase ended before ada answered. The passages
it rests on (P lines 329--331, 1030--1034) were reread for this note and
say what \#26 quotes, and Section~\ref{sec:table2} uses the corrected
version.
\item \textbf{Full versus partial implementation} (ada \#17(4), \#18(1),
on emmy \#14(3) and \#15). Resolved: emmy accepted it (\#21(1)), and the
verifier raised it independently (\#20).
\item \textbf{The format-review step} (emmy \#26(1)). P describes it only
for AR (P line 320). If Single-reviewer output skips it, AR's comment
count rises mechanically. Open.
\item \textbf{Unmatched does not mean false} (\#21(4), \#26(5)). P's
judge scores similarity to human comments (Cohen's $\kappa=0.75$ against
annotators, P line 272) and may penalise valid comments (P lines
685--688). This limits every reading of Table~2, and it is why falsity
in the proposed test is confirmed by a person. Stands.
\item \textbf{Micro versus macro F1} (\#26). The ratios in \#26(3) assume
micro-averaging, and P does not say which it uses. Open.
\item \textbf{P's LCB ladder confounds interaction with the acceptance
rule} (ada \#9). Single-reviewer and Two-reviewers edit once and MARS is
capped at two rounds, while AR re-edits until R and C converge at once
and find no flaws (P lines 203--205). P states no outer cap for LCB, and
its SWE-bench workflow caps the outer loop at two (P lines 1884, 1892).
The LCB gain may therefore come from the stricter acceptance rule and the
extra edits rather than from R-C interaction. emmy confirmed this from
P's text and found that P's generic AR skill allows seven inner rounds,
against five in P Sec.~3.5 (\#13, P lines 199, 1689); ada confirmed it
(\#17(1)). Open: which caps ran on LCB, and the separating run.
\item \textbf{Sign versus co-movement} (emmy \#4, correcting \#2 and \#3).
Resolved; ada checked it (\#11).
\item \textbf{Pure pooling in Q's example} (emmy \#13, correcting \#8).
Resolved (\#17(1)).
\item \textbf{Report matching for co-moving biases} (\#5(a)). Withdrawn
by emmy (\#14(4)); message mechanisms for $k\ne1$ remain open (\#15(4)).
\item \textbf{P's evidence is thin} (\#5(c), \#11(2), \#26(5)). About five
of 105 LCB tasks separate MARS ($82\%$) from AR ($87\%$). SWE-PRBench has
100 PRs and no intervals, and the text constraint was designed on the
benchmark on which it is scored. Stands.
\item \textbf{The scalar model is an analogy} (\#5(b), \#11(3), \#15(2),
\#24(3)). P's verdicts are discrete, P reports no per-task signed errors,
and the model cannot express failures of discrimination such as Case~B's
(\#26(6)). Stands.
\end{itemize}

\section{Sources from outside Q and P}

None of these enters a derivation. They position the results, bear on
the premise, or supply inputs for the proposed test.
\begin{itemize}
\item Krishna and Morgan (2001), \emph{A Model of Expertise}, Quarterly
Journal of Economics 116(2), 747--775 (\#3). The like-versus-opposing
experts point is theirs. emmy argues that co-movement is a different
condition from their sign condition (\#4).
\item Moore and Repullo (1988) and Dutta and Sen (1991), on
implementation with two agents (\#14(4)).
\item Correlated-agreement peer prediction: Dasgupta and Ghosh (2013),
Shnayder et al.\ (2016), and the survey arXiv 2106.03176 (\#8). They
motivated the withdrawn audit.
\item LLM conformity: arXiv 2508.18321 and arXiv 2606.01637 (\#8).
\item Correlated errors across LLMs: Kim et al., arXiv 2506.07962 (\#3,
\#23(2)), and arXiv 2605.29800 (\#3).
\item Self-critique against cross-model critique: arXiv 2606.05976
(\#21 from a search snippet, \#23(2) with figures), arXiv 2505.17656
(\#21, snippet only, not read), DeltaBench (arXiv 2502.19361), arXiv
2408.10495, and arXiv 2605.21537, which has no cross-model arm and so
cannot separate ownership from disposition (\#23(2)).
\item SWE-PRBench: the public dataset (CC BY 4.0), with diffs and human
comments but no model reviews, and its paper, arXiv 2603.26130, which
reports single-pass hallucination rates of $0.19$ to $0.42$ on diff-only
input ($0.227$ for Sonnet~4.6), so own-model false flags should exist on
most diffs (\#23(1), \#23(3)).
\item Morris and Shin (2002) and Geanakoplos and Polemarchakis, for the
information version that was not worked out (\#9, ada's notes).
\end{itemize}
Searches are recorded in \#3, \#8, \#21 and \#23 (\#24(5)), and none was
run for the contrarian result (\#15(6)). These are search records, not
novelty claims. A crossover of critic and error source is not new as a
method (\#21(7)).

\section{What remains open}

The material is thin: a stylised quadratic model with complete
information between the agents, a reading of P's text and prompts, and
P's published aggregates. Nothing about P has been measured.
\begin{enumerate}
\item \emph{The premise on P's side.} Whether P's critic reacts to false
flags its own model would write, which under the hard reading is whether
truth is focal for a same-model R-C pair (\#20, \#21(3), \#24). Every
claim about P depends on it.
\item \emph{What carries P's Table~2.} Case~A additions or R's deletions
for naive AR's deficit (\#26(3)); union plus no deletion, or
evidence-grounded disagreement, for v2's gain (\#26(4)); whether the
format-review step is applied to AR alone, and micro or macro F1
(\#26(1)). These need P's logs, which are not public (P line 728), or
reruns.
\item \emph{LCB.} Which inner cap ran, and whether an outer cap applied
(\#13, \#24(2)); Single-reviewer inside AR's outer loop (\#9).
\item \emph{Predictions not run.} If the crossover contrast of
Section~\ref{sec:next} is positive, a cross-family critic cuts Case~A
over-flagging (\#21(6)), and AR's gain should grow as R's and C's signed
errors become negatively correlated (\#3, \#4). Extra rounds do not help
same-direction clones (\#3). Under the text constraint, residual
fabricated comments come from C's certificate-free additions, and the
gain falls on flags settled by citations that occur in the diff (\#3,
\#10). A one-round contrarian design with an exaggerating critic beats
P's R on over-flagging (\#15(1), \#18(2)).
\item \emph{Theory.} Message mechanisms for $k\ne1$ (\#14(4), \#15(4));
difference rewards for R that are neither concave nor convex along the
path (\#15(5)); non-quadratic contrarian rewards (\#18(3)); the minimax
case above $1$, computed by hand only (\#18(2)); the information version,
in which a beauty contest gives evidence weight
$w=(1-r)\lambda/(1-r\lambda)<\lambda$ but P's sequential loop calls for a
model of conformity garbling C's message (\#9, \#11(4)); and failures of
discrimination, which the scalar model cannot express (\#26(6)).
\item \emph{Novelty} is unestablished (\#5(d), \#11(5), \#15(6),
\#21(7), \#24(5)).
\item \emph{Coverage.} The peer phase ended when the allowance ran out.
emmy's last ready is at \#26, but ada's is at \#24 and does not cover
\#26, so emmy's reading of Table~2 has only emmy's own check and the
rereading of P's passages for this note. The revelation-principle reading
(\#10(3)) was flagged by emmy as unchecked (\#5(e)) and not taken up
again.
\end{enumerate}

\section{The smallest next step}\label{sec:next}

Every claim about P turns on one quantity: whether P's critic, on the
same diff, reacts to a false flag that its own model would write (\#20,
\#21(3)). The smallest step that settles it is the core of emmy's
same-diff test (\#21(4)), with two arms. On each of P's 100 SWE-PRBench
diffs, run P's baseline critic prompt (P lines 1083--1120) on Sonnet~4.5,
in a fresh context, against R's round-1 review $r_i$ (arm A0) and against
$r_i$ with one appended false flag $f_i^{\mathrm{own}}$ (arm A1). The flag
is one that Sonnet~4.5 wrote as R on that diff in an extra sample; it
cites real lines of the diff, and a person has confirmed it false, since
P's judge cannot certify falsity. It is appended rather than swapped for
a true flag, so that C's search for missed bugs cannot react to an
omission instead. Report the share of diffs on which C names the injected
flag as spurious, which is the primary outcome, and the paired change in
AGREE between the arms with its interval, whose $95\%$ half-width at
$n=100$ is $0.085$ to $0.126$ (\#21(5)). The inputs are public except R's
reviews, which one R pass per diff supplies (\#23(1)). The cost is some
hundreds of model calls and a person checking one flag per diff.

If C seldom names the flag and AGREE barely moves, C's AGREE does not
react to its own model's errors, and P's naive loop is a consensus device
for common-mode errors under either reading of the loop (\#21(6)). P's v2
gain then has to come from the citation channel (\#10) or from union plus
no deletion (\#26(4)). If C names most such flags, the premise fails for
this pair, and the naive-AR deficit has to come from something other than
co-moving errors. \#21(6) named Case~B at this point; by \#26(2) Case~B is
neutral against Single-reviewer, which leaves the deletion channels of
\#26(3).

The core does not say whether a low detection rate reflects co-movement
or weak discrimination of false flags in general. That needs the
other-family arm and the crossover contrast
$I=[d(S,G)-d(S,S)]+[d(G,S)-d(G,G)]$ of \#21(4), where $d(c,s)$ is critic
$c$'s detection rate for false flags from source $s$, $S$ is Sonnet and
$G$ another family. Its $95\%$ half-width at $n=100$ is about $0.18$ to
$0.20$ unless several flags are injected per diff (\#21(5)). Separating
ownership from disposition needs ada's arm that shows the same flag to C
as its own earlier draft (\#23(3)). Before any spending, P's authors can
be asked whether Single-reviewer passes through the format-review step
and whether F1 is micro- or macro-averaged (\#26(1)), which caps ran on
LCB (\#13), and for the per-comment judge labels, which split Table~2 into
precision, recall and comments per PR at no call cost (\#23(1), \#26(3)).
Those answers would settle the confounds in P's comparisons, but they do
not measure C's response to its own model's errors.

\end{document}
